% ~10 min
\partframe{The Laws}{Divide and Conquer as Algebra}

\begin{frame}
  \frametitle{One Shape for Many Algorithms}
  \begin{center}
    \begin{tikzpicture}[alt={The shape of a divide-and-conquer algorithm:
        split the input into smaller pieces, solve every piece recursively,
        then combine the partial results into the answer.},
      stage/.style = {box, minimum height = 10mm, text width = 22mm,
        align = center}]
      \node[stage] (split) {\textbf{split}\\ $c$};
      \node[stage, right = 16mm of split] (rec) {solve each piece\\
        (recursion)};
      \node[stage, right = 16mm of rec] (comb) {\textbf{combine}\\ $a$};
      \node[left = 6mm of split] (in) {input};
      \node[right = 6mm of comb] (out) {output};
      \draw[->, thick] (in) -- (split);
      \draw[->, thick] (split) -- node[above, note] {pieces} (rec);
      \draw[->, thick] (rec) -- node[above, note] {results} (comb);
      \draw[->, thick] (comb) -- (out);
    \end{tikzpicture}
  \end{center}
  \begin{center}
    \renewcommand\arraystretch{1.3}
    \begin{tabular}{lll}
      \textbf{Algorithm} & \textbf{split} $c$ & \textbf{combine} $a$ \\
      \hline
      \uncover<2->{quicksort} & \uncover<2->{partition around a pivot} &
        \uncover<2->{concatenate} \\
      \uncover<3->{mergesort} & \uncover<3->{cut in half} &
        \uncover<3->{merge} \\
      \uncover<4->{$x^n$} & \uncover<4->{$n \mapsto$ last bit, $n/2$} &
        \uncover<4->{square, maybe multiply by $x$} \\
    \end{tabular}
  \end{center}
  \uncover<5->{\textbf{Hylomorphism} $\hylo{a}{c}$: the recursion is
    \emph{generic}; only $a$ and $c$ change.}
  \speakernote{The key abstraction. Many algorithms have the same
    skeleton: split the input into smaller pieces (a coalgebra, c), solve
    the pieces recursively, combine the results (an algebra, a). [click]
    Quicksort: split = partition, combine = concatenate. [click] Mergesort:
    the work moves from split to combine. [click] Fast exponentiation: split
    a number into its last bit and the rest. [click] The recursion itself
    is written once, as a higher-order function of a and c. This is a
    hylomorphism (Greek: matter and form). Laws about hylomorphisms then
    hold for every algorithm of this shape.}
\end{frame}

\begin{frame}
  \frametitle{The Call Tree of Quicksort}
  \begin{center}
    \begin{tikzpicture}[alt={The call tree of quicksort on the list 3, 1,
        4, 5, 2. Splitting around pivot 3 gives 1, 2 and 4, 5, which are
        split again around pivots 1 and 4. Combining concatenates the
        results from the leaves up, giving 1, 2, 3, 4, 5.},
      level distance = 13mm,
      level 1/.style = {sibling distance = 46mm},
      level 2/.style = {sibling distance = 22mm},
      every node/.style = {box}]
      \node (r) {$[3,1,4,5,2]$}
        child {node (l) {$[1,2]$}
          child {node {$[\,]$}}
          child {node {$[2]$}}}
        child {node (rr) {$[4,5]$}
          child {node {$[\,]$}}
          child {node {$[5]$}}};
      \node[hot, right = 4mm of r] {pivot 3};
      \node[hot, right = 4mm of l] {pivot 1};
      \node[hot, right = 4mm of rr] {pivot 4};
      \draw[->, very thick, accent] (-6.0, 0.2) -- node[left, draw = none,
        fill = none, note, text = accent] {$c$: split} (-6.0, -2.6);
      \draw[->, very thick, accent2] (6.0, -2.6) -- node[right, draw = none,
        fill = none, note] {$a$: combine} (6.0, 0.2);
    \end{tikzpicture}
  \end{center}
  \uncover<2->{The tree is \emph{virtual}: it is never built in memory.}
  \speakernote{Splitting goes down, combining comes up. Conceptually the
    algorithm builds this tree, then folds it. [click] But the hylomorphism
    consumes every node as soon as it is produced, so the tree only exists
    as the call stack. Whether an intermediate structure exists or not is
    exactly what the laws let us control.}
\end{frame}

\begin{frame}
  \frametitle{Producers and Consumers}
  \begin{center}
    \begin{tikzpicture}[alt={Top: an unfold produces a tree from the input
        using the split step c, and a fold consumes the tree using the
        combine step a. Bottom: the hylomorphism goes straight from input to
        output, with no tree in between.},
      stage/.style = {box, minimum height = 9mm, minimum width = 24mm}]
      \node (in) {input};
      \node[stage, right = 8mm of in] (un) {\textbf{unfold} with $c$};
      \node[hot, right = 8mm of un] (t) {tree};
      \node[stage, right = 8mm of t] (fo) {\textbf{fold} with $a$};
      \node[right = 8mm of fo] (out) {output};
      \draw[->, thick] (in) -- (un);
      \draw[->, thick] (un) -- (t);
      \draw[->, thick] (t) -- (fo);
      \draw[->, thick] (fo) -- (out);
      \node[box, minimum height = 9mm, below = 12mm of t,
        minimum width = 60mm, draw = good, fill = good!8] (h)
        {\textbf{hylo} $\hylo{a}{c}$: no tree};
      \draw[->, thick] (in) |- (h);
      \draw[->, thick] (h) -| (out);
    \end{tikzpicture}
  \end{center}
  \begin{itemize}[<+->]
    \item Specification: \emph{build}, then \emph{consume}.
      Easy to read.
    \item Program: do both at once. Fast.
  \end{itemize}
  \speakernote{An unfold produces a data structure (here the call tree);
    a fold consumes one. Specifications are naturally written as
    build-then-consume pipelines, like the sort-then-double example: each
    stage is simple. Programs should not materialise the intermediate
    structure: that is allocation, memory traffic, and garbage. A
    hylomorphism is the fused version. Removing the intermediate structure
    is called deforestation (Wadler, 1990).}
\end{frame}

\begin{frame}
  \frametitle{Two Laws}
  \begin{center}
    \begin{tikzpicture}[alt={Fusion: a hylomorphism followed by a
        post-processing step h equals a single hylomorphism whose combine
        step does the post-processing. Deforestation: a hylomorphism that
        builds a structure, followed by one that takes it apart, equals one
        hylomorphism.},
      stage/.style = {box, minimum height = 10mm, minimum width = 14mm,
        font = \large},
      lbl/.style = {font = \bfseries, anchor = east}]
      \node[lbl] at (-0.4, 0) {Fusion};
      \node[stage, anchor = west] (h1) at (0, 0) {$\hylo{a}{c}$};
      \node[stage, right = 5mm of h1] (p) {$h$};
      \draw[->, thick] (h1) -- (p);
      \node[right = 5mm of p, font = \Large] (e1) {$=$};
      \node[stage, right = 5mm of e1, draw = good, fill = good!8]
        {$\hylo{b}{c}$};
      \node[note] at ([yshift = -9mm]e1) {if $h$ distributes over $a$};
      \node[lbl] at (-0.4, -2.2) {Deforestation};
      \node[stage, anchor = west] (h2) at (0, -2.2) {$\hylo{\op{build}}{c}$};
      \node[stage, right = 5mm of h2] (q) {$\hylo{a}{\op{take\_apart}}$};
      \draw[->, thick] (h2) -- (q);
      \node[right = 5mm of q, font = \Large] (e2) {$=$};
      \node[stage, right = 5mm of e2, draw = good, fill = good!8]
        {$\hylo{a}{c}$};
    \end{tikzpicture}
  \end{center}
  \begin{itemize}[<+->]
    \item Proved \textbf{once}, for every split, combine, data type.
    \item Each use is a \emph{rewrite step}.
  \end{itemize}
  \speakernote{Fusion: if a post-processing function h commutes with the
    combine step (h after a equals b after applying h to the pieces), then
    the post-processing can be pushed into the combine step: one pass
    instead of two. Deforestation: building a structure and immediately
    taking it apart cancel out. Both follow from one fact: a hylomorphism
    with a terminating split has exactly one solution. [click] These hold
    generically, for any data type with a well-behaved split. [click] An
    optimisation is a sequence of such steps.}
\end{frame}

\begin{frame}
  \frametitle{Deriving the Example}
  \begin{center}
    \begin{tikzpicture}[alt={Three steps. First, sort then double. Second,
        sort is a hylomorphism with split partition and combine
        concatenate, followed by double. Third, by fusion, doubling moves
        into the combine step, giving a single hylomorphism that
        concatenates and doubles the pivot.},
      stage/.style = {box, minimum height = 8mm},
      eq/.style = {font = \Large, anchor = east}]
      \node[stage, anchor = west] (a1) at (0, 0) {sort};
      \node[stage, right = 5mm of a1] (a2) {double};
      \draw[->, thick] (a1) -- (a2);
      \uncover<2->{
        \node[eq] at (-0.2, -1.4) {$=$};
        \node[stage, anchor = west] (b1) at (0, -1.4)
          {$\hylo{\op{concat}}{\op{partition}}$};
        \node[stage, right = 5mm of b1] (b2) {double};
        \draw[->, thick] (b1) -- (b2);
        \node[note, right = 5mm of b2] {definition of quicksort};
      }
      \uncover<3->{
        \node[eq] at (-0.2, -2.8) {$=$};
        \node[stage, anchor = west, draw = good, fill = good!8] (c1)
          at (0, -2.8)
          {$\hylo{\op{concat} \text{ and double pivot}}{\op{partition}}$};
        \node[note, right = 5mm of c1] {fusion};
      }
    \end{tikzpicture}
  \end{center}
  \uncover<4->{Each $=$ is a \textbf{theorem} applied by a proof assistant.}
  \speakernote{The opening example in this vocabulary. Quicksort is a
    hylomorphism: split = partition, combine = concatenate. [click]
    Doubling distributes over concatenation (doubling a concatenation is
    concatenating the doubled parts), so fusion applies. [click] The result
    doubles each pivot as it is placed: the hand-written program from the
    start. [click] In our work each of these steps is a rewrite in Rocq,
    checked by its kernel. The rest of the talk: what it takes to make that
    work, and what comes out.}
\end{frame}
